
@@M0@@

Indiquons pour finir l'origine des formules (2.3.1) et (2.3.3). On considère l'exemple 2.2.b), où les opérations @@M1@@ sont définies par les puissances extérieures des @@M2@@-Modules. On souhaite que @@M3@@ soit un @@M4@@-homomorphisme~; on doit donc définir sur @@M5@@ des lois telles que pour tous @@M6@@-Modules @@M7@@, on ait~:
@@M8@@
or les relations (2.3.1) et (2.3.3) sont précisément les relations précédentes appliquées au cas où @@M9@@ et @@M10@@ sont des modules inversibles.

Nous verrons dans Exp. VI dans quelle mesure la propriété souhaitée est vérifiée.

\textbf{Définition 2.4.} \uline{Un pré}-@@M11@@-\uline{anneau} @@M12@@ \uline{sera appelé un} @@M13@@-\uline{anneau si l'homomorphisme de groupes abéliens} @@M14@@ \uline{est un} @@M15@@-\uline{homomorphisme} \uline{pour la structure de pré}-@@M16@@-\uline{anneau sur le deuxième membre définie en} 2.3.

Cela signifie donc qu'on a les relations~: @@M17@@,
@@M18@@

En termes des applications @@M19@@, les deux dernières relations équivalent à : @@M20@@, @@M21@@,
@@M22@@

